\documentclass[10pt,a4paper]{amsart}
\usepackage[margin=19mm]{geometry}
\usepackage{amsmath,amssymb,mathtools}
\usepackage[scr=rsfs,cal=euler]{mathalfa}
\usepackage{xfrac}
\usepackage[unicode,hidelinks,bookmarks=false]{hyperref}
\newcommand{\ie}{i.e.\ }
\newcommand{\cf}{cf.\ }
\newcommand{\resp}{respectively\ }
\newcommand{\Subsection}{\subsection}
\providecommand{\nobreakdash}{\nobreak\hbox{-}}
\setlength{\emergencystretch}{2em}
\multlinegap0pt
\usepackage{bm}
\usepackage{bbm}
\usepackage{dsfont}
\usepackage{xspace}
\usepackage{cancel}
\usepackage{mathtools}
\usepackage[shortlabels]{enumitem}
\usepackage{amsthm}
\makeatletter
\@ifundefined{theorem}{\newtheorem{theorem}{Theorem}}{}
\@ifundefined{proposition}{\newtheorem{proposition}[theorem]{Proposition}}{}
\makeatother
\newcommand{\R}{\mathbb{R}}
\newcommand{\ee}{\bm e}
\newcommand{\jj}{\bm j}
\newcommand{\kk}{\bm k}
\newcommand{\one}{\mathbf{1}}
\newcommand{\qq}{\bm q}
\title[Structure-Preserving Numerical Schemes for Two-Stage Local a]{Structure-Preserving Numerical Schemes for Two-Stage Local and Nonlocal Dispersal Systems}
\author{Markjoe O. Uba}
\date{Source: arXiv:2609.29524v1 (CC BY 4.0)}
\begin{document}
\maketitle

\noindent\emph{Source and licence.} Structure-Preserving Numerical Schemes for Two-Stage Local and Nonlocal Dispersal Systems. Version arXiv:2609.29524v1 (CC BY 4.0), distributed under the Creative Commons Attribution 4.0 International licence, as recorded in the arXiv licence metadata for this exact version and shown on the abstract page https://arxiv.org/abs/2609.29524v1. Full rights notice: licenses/arxiv-2609.29524-CC-BY-4.0.txt.\par\smallskip

\noindent\textit{Editorial scope.} Counted unit: the Proposition whose source label is \texttt{prop:matrix-structure} in the article (source0.tex lines 789--807, source label \texttt{prop:matrix-structure}), delivered together with the article's own complete original proof of it (source0.tex lines 809--865, one \texttt{proof} environment of 57 lines). No mathematics is written, repaired or re-derived: statement and proof are the article's own text line for line. Required context retained uncounted: eq:J-positive-ball (lines 678--682); eq:Dh (lines 767--775); eq:local-Dh (lines 777--784). Every definition, display and label of those lines is retained unchanged. Omitted from this package and not redistributed: 1--677, 683--766, 776--776, 785--788, 808--808, 866--2132. Nothing inside a retained excerpt is omitted or altered: every retained body is delivered exactly as the article's lines are written, so no omission receipt of any kind accompanies this package. Citations: the retained excerpts carry no citation command at all, so no bibliography entry of the article is transcribed and no reference is redistributed. Reference adaptation: none was needed and none was made; every reference inside the delivered unit resolves inside it, so no unresolved reference or citation is produced. Printed slips kept literal: no printed word, symbol, hypothesis or number of the counted statement or of its proof was altered. Layout and class: the article's own class is \texttt{article} with packages including geometry, amsmath, amssymb, amsthm, mathtools, bm, booktabs, array, graphicx, microtype, enumitem, hyperref, lineno, xcolor; the wrapper is the lane's pinned \texttt{amsart} editorial wrapper at 10pt on A4, which restates the theorem-like environments the retained text needs and adds the article's own macro definitions that the retained text needs (\texttt{\textbackslash R}, \texttt{\textbackslash ee}, \texttt{\textbackslash jj}, \texttt{\textbackslash kk}, \texttt{\textbackslash one}, \texttt{\textbackslash qq}), with extra wrapper packages (bm, bbm, dsfont, xspace, cancel, mathtools, enumitem). The delivered numbering is the unit's own. Assets: no retained span contains a figure environment, a graphics-inclusion command, a PDF-page inclusion, a table or a TikZ picture; no image file is copied into this package, the package declares no asset dependency and no image file is redistributed with it. Licence: version arXiv:2609.29524v1 of this article is distributed under the Creative Commons Attribution 4.0 International licence, as recorded in the arXiv licence metadata for this exact version and shown on the abstract page https://arxiv.org/abs/2609.29524v1. A full rights notice is delivered at licenses/arxiv-2609.29524-CC-BY-4.0.txt. Characters: every retained excerpt is pure ASCII, so the delivered engine, XeLaTeX with its default Latin Modern text font, reports no missing character.
\begin{equation}
\label{eq:J-positive-ball}
 J(z)>0
 \qquad\text{for }|z|\le\rho_J.
\end{equation}

\begin{equation}
\label{eq:Dh}
 (D_{\delta,h}^{\Box}V)_{\jj}
 =
 \frac{2}{\kappa_{\delta,h}}
 \sum_{\qq\in\mathcal Q_{\delta,h}}
 \omega_{\qq}^{\delta,h}
 \big(V^E_{\jj+\qq}-V_{\jj}\big).
\end{equation}

\begin{equation}
\label{eq:local-Dh}
 (D_{0,h}^{\Box}V)_{\jj}
 =
 \sum_{\ell=1}^{d}
 \frac{V^E_{\jj+\ee_{\ell}}-2V_{\jj}
       +V^E_{\jj-\ee_{\ell}}}{h^2}.
\end{equation}

\begin{proposition}[Structural properties of the reflected Cartesian operator]
\label{prop:matrix-structure}
For $0\le\delta\le\delta_{\Box}$:
\begin{enumerate}[label=(\roman*)]
\item $(D_{\delta,h}^{\Box})_{ij}\ge0$ for $i\ne j$, and $D_{\delta,h}^{\Box}\one=0$;
\item $D_{\delta,h}^{\Box}$ is symmetric and negative semidefinite;
\item for every $V\in\R^{N_h^{\Box}}$,
\begin{equation}
\label{eq:energy}
 -h^dV^TD_{\delta,h}^{\Box}V
 =
 \frac{h^d}{2}
 \sum_{\substack{i,j=1\\ i\ne j}}^{N_h^{\Box}}
 w_{ij}^{\delta,h}(V_i-V_j)^2
 \ge0;
\end{equation}
\item $\ker D_{\delta,h}^{\Box}=\operatorname{span}\{\one\}$.
\end{enumerate}
\end{proposition}

\begin{proof}
Let $R:\mathbb Z^d\to\mathcal I_h$ denote the coordinatewise reflection map
defined by $V^E_{\kk}=V_{R(\kk)}$.  In one coordinate,
\[
 R_\ell^{-1}(k)
 =
 \{\,k+2mN_\ell,\,-k-1+2mN_\ell:m\in\mathbb Z\,\}.
\]
Thus, every reflected contribution from one cell to another has a
corresponding reverse contribution generated by an offset of the same Euclidean
length.  Since $J$ is radial, the two contributions have the same weight.
Consequently,
\[
w_{ij}^{\delta,h}=w_{ji}^{\delta,h}.
\]

For $i\ne j$, the coefficients in \eqref{eq:Dh} therefore satisfy
\[
 w_{ij}^{\delta,h}\ge0,\qquad w_{ij}^{\delta,h}=w_{ji}^{\delta,h},\qquad
 (D_{\delta,h}^{\Box})_{ii}=-\sum_{j\ne i}w_{ij}^{\delta,h}.
\]
Consequently,
\[
 (D_{\delta,h}^{\Box}\one)_i=-\sum_{j\ne i}w_{ij}^{\delta,h}+\sum_{j\ne i}w_{ij}^{\delta,h}=0,
 \qquad
 (D_{\delta,h}^{\Box})^T=D_{\delta,h}^{\Box}.
\]
Therefore
\[
 -h^dV^TD_{\delta,h}^{\Box}V
 =
 \frac{h^d}{2}
 \sum_{\substack{i,j=1\\ i\ne j}}^{N_h^{\Box}}
 w_{ij}^{\delta,h}(V_i-V_j)^2
 \ge0.
\]

If the nonlocal formula \eqref{eq:Dh} is used, then \eqref{eq:J-positive-ball} gives $\omega_{\pm\ee_\ell}^{\delta,h}>0$ for $\ell=1,\ldots,d$.  Thus equality in \eqref{eq:energy} implies equality of the values at every
pair of neighboring cell centers in each coordinate direction.  Repeating
these equalities along the Cartesian coordinate lines gives $V_{\jj}=V_{\kk}$ for all $\jj,\kk\in\mathcal I_h$.  If the local formula \eqref{eq:local-Dh} is used, then
\[
 -h^dV^TD_{0,h}^{\Box}V
 =
 h^{d-2}
 \sum_{\ell=1}^d
 \sum_{\substack{\jj\in\mathcal I_h\\0\le j_\ell\le N_\ell-2}}
 \big(V_{\jj+\ee_\ell}-V_{\jj}\big)^2.
\]
Hence
\[
 V^TD_{0,h}^{\Box}V=0
 \quad\Longrightarrow\quad
 V_{\jj+\ee_\ell}=V_{\jj}
\]
for every $\ell=1,\ldots,d$ and every
$\jj\in\mathcal I_h$ with $0\le j_\ell\le N_\ell-2$.  Successive coordinate steps therefore give $V_{\jj}=V_{\kk}$ for all $\jj,\kk\in\mathcal I_h$.  Thus in both cases $\ker D_{\delta,h}^{\Box}=\operatorname{span}\{\one\}$.
\end{proof}

\end{document}
